% ======================================================================
\section{Single inclusive gluon production: setting up the cross section}\label{sec:cross}
% ======================================================================

We now use the operator $\Omega$ to write the single inclusive gluon spectrum in the collision of the dilute projectile with a dense target. In this section the spectrum is expressed through the coefficients of $\Omega$ and the Wilson lines of the target, to order $g^4$; the coefficients are not yet inserted. The result is organised in three groups of terms, which are the starting point of the evaluation in the following sections.

\subsection{The observable}\label{sec:observable_def}

\paragraph{Eikonal scattering.} Before the collision the projectile is in the state $\Omega\ket0$, i.e.\ $\Omega\ket0\otimes\ket v$ with the valence state $\ket v$. The target is described by the eikonal $S$-matrix~\cite{Kovner:2005en,Lublinsky:2016meo}: at high energy a soft gluon and the valence charges keep their transverse positions and longitudinal momenta while they cross the target, and their colour is rotated by the adjoint Wilson line $U(\v x)$ of the target field at their transverse position,
\begin{align}\label{eq:Srotation}
S^\dagger\rho^a(\v x)\,S&=U^{ab}(\v x)\,\rho^b(\v x)\,,\notag\\
S^\dagger a^a_i(k^+,\v x)\,S&=U^{ab}(\v x)\,a^b_i(k^+,\v x)\,,\qquad
S^\dagger a^{\dagger a}_i(k^+,\v x)\,S=U^{ab}(\v x)\,a^{\dagger b}_i(k^+,\v x)\,.
\end{align}
The matrix $U$ is real and orthogonal, $U^\dagger=U^T$, $U^TU=1$, and $S\ket0=\ket0$ for the soft vacuum. After the collision the state is $S\,\Omega\ket0$.

\paragraph{The number of produced gluons.} A gluon observed in the final state is a dressed gluon, i.e.\ a state $\Omega\,a^\dagger\ket0$ and not a bare quantum of the cloud of the projectile. The number of dressed gluons is counted by the number operator rotated by $\Omega$, $\Omega\,a^\dagger a\,\Omega^\dagger$~\cite{Kovner:2006wr}. The single inclusive spectrum is therefore
\begin{align}\label{eq:observable}
\frac{d^3N}{d^3k}=\frac{1}{(2\pi)^3}\int_{w,w'}e^{-i\v k\cdot(\v w-\v w')}\,
\Big\langle\bra0\,\Omega^\dagger S^\dagger\Omega\;a^{\dagger a}_i(k^+,\v w')\,a^a_i(k^+,\v w)\;\Omega^\dagger S\,\Omega\ket0\Big\rangle\,,
\end{align}
where $\langle\dots\rangle$ is the average over the valence state, taken last. The factor $(2\pi)^{-3}$ and the phase follow from Eqs.~\eqref{eq:aadagger} and~\eqref{eq:FTconventions}: the density of gluons in $d^3k$ is $(2\pi)^{-3}a^{\dagger a}_i(k)a^a_i(k)$, and $a^{\dagger}(k^+,\v k)\,a(k^+,\v k)=\int_{w,w'}e^{-i\v k\cdot(\v w-\v w')}a^{\dagger}(k^+,\v w')\,a(k^+,\v w)$. The amplitude stands at the unprimed point $\v w$ and the conjugate amplitude at the primed point $\v w'$. We use this convention, the phase $e^{-i\v k\cdot(\text{unprimed}-\text{primed})}$ with the conjugated factor to the left, in all expressions below. Without a target, $S=1$, the spectrum vanishes to all orders because $\Omega^\dagger\Omega=1$.

\paragraph{Barred coefficients.} We insert $SS^\dagger=1$ to the left of the number operator and between $\Omega^\dagger$ and $S$, and define
\begin{align}\label{eq:Omegabar}
\bar\Omega\equiv S^\dagger\,\Omega\,S\,.
\end{align}
By Eq.~\eqref{eq:Srotation}, $\bar\Omega$ is given by the expansion~\eqref{eq:OmegaX} with every charge and every gluon operator rotated. A bar on a coefficient denotes the replacement of every charge in it by the rotated charge,
\begin{align}\label{eq:bar}
\bar{\mf A}\equiv\mf A\big|_{\rho^a(\v u)\to U^{ab}(\v u)\rho^b(\v u)}\,,\qquad\text{and the same for }\ \bar{\mf N},\ \bar{\mf B}_{1,2,3}\,.
\end{align}
The coefficients $\mf C_{1,2}$ do not depend on $\rho$, and $\bar{\mf C}_{1}=\mf C_1$; we keep the bar as a reminder that the coefficient originates from $\bar\Omega$. Since the barred charges obey the same algebra~\eqref{eq:rhoalgebra}, the barred coefficients obey the unitarity relations of Section~\ref{sec:unitarity}. With $S^\dagger a^\dagger a\,S=Ua^\dagger\,Ua$ the spectrum becomes
\begin{align}\label{eq:observable2}
\frac{d^3N}{d^3k}=\frac{1}{(2\pi)^3}\int_{w,w'}e^{-i\v k\cdot(\v w-\v w')}\,\Big\langle\bra0\,\Omega^\dagger\bar\Omega\;U^{ab'}(\v w')\,a^{\dagger b'}_i(k^+,\v w')\;U^{ab}(\v w)\,a^b_i(k^+,\v w)\;\bar\Omega^\dagger\Omega\ket0\Big\rangle\,.
\end{align}

\paragraph{Amplitudes.} Between the two gluon operators we insert the complete set of soft Fock states. The state $a(k^+,\v w)\,\bar\Omega^\dagger\Omega\ket0$ is projected on the vacuum, on one-gluon states, and so on; these projections are the overlaps of $\bar\Omega^\dagger\Omega\ket0$ with states with one, two, \dots\ gluons, one of which is the measured gluon. With the normalisation~\eqref{eq:states},
\begin{align}\label{eq:N_amplitudes}
\frac{d^3N}{d^3k}=\frac{1}{(2\pi)^3}\Big\langle\mathcal A^{a\dagger}_i(k)\,\mathcal A^a_i(k)+\int dq^+\!\int_{v}\mathcal A^{ab\dagger}_{ij}(k;q^+,\v v)\,\mathcal A^{ab}_{ij}(k;q^+,\v v)\Big\rangle+\dots,
\end{align}
with the one-gluon and two-gluon amplitudes
\begin{align}\label{eq:amplitudes_def}
\mathcal A^a_i(k)&=\sqrt{2\pi}\int_{w}e^{-i\v k\cdot\v w}\,U^{ab}(\v w)\,\bra{g^b_i(k^+,\v w)}\bar\Omega^\dagger\Omega\ket0\,,\notag\\
\mathcal A^{ab}_{ij}(k;q^+,\v v)&=\sqrt{2\pi}\int_{w}e^{-i\v k\cdot\v w}\,U^{ac}(\v w)\,\bra{g^b_j(q^+,\v v)\,g^c_i(k^+,\v w)}\bar\Omega^\dagger\Omega\ket0\,.
\end{align}
The amplitudes are operators on $\mathcal H_{\rm val}$. The second gluon of $\mathcal A^{ab}_{ij}$, with longitudinal momentum $q^+$ at the transverse position $\v v$, is not measured and is integrated over. The dots in Eq.~\eqref{eq:N_amplitudes} are the contributions of final states with three or more gluons; the amplitude for $n$ gluons is of order $g^n$, so that they are of order $g^6$. The amplitudes are evaluated with the expansion~\eqref{eq:OmegaX}: $\Omega$ acting on $\ket0$ involves only $\mf N$, $\mf A_2$ and $\mf B_2$, Eq.~\eqref{eq:Omegavac}, and $\bar\Omega^\dagger$ then acts on states with up to two gluons. This is where the coefficients $\mf B_3$ and $\mf C_1$ of the dressed one-gluon state enter, and it confirms the counting of Section~\ref{sec:threeout}: the one-gluon amplitude is needed to order $g^3$ and the two-gluon amplitude to order $g^2$, for which the coefficients of Table~\ref{tab:coeffs} are sufficient.

From here on $\mf A\equiv\mf A^{(1)}_2$, $\mf A^{(3)}\equiv\mf A^{(3)}_2$, and the coefficients $\mf A_1$, $\mf B_1$, $\mf C_2$ of $\bar\Omega^\dagger$ are eliminated with the unitarity relations~\eqref{eq:unitX}.

\subsection{Expansion in the coupling and the leading order}\label{sec:LO}

The amplitudes are expanded in $g$,
\begin{align}\label{eq:ampexp}
\mathcal A^a_i(k)=\mathcal A^{a}_{{\rm LO}\,i}(k)+\mathcal A^a_{(3)\,i}(k)+O(g^5)\,,\qquad \mathcal A^{ab}_{ij}=O(g^2)\,,
\end{align}
with $\mathcal A_{\rm LO}=O(g)$ and $\mathcal A_{(3)}=O(g^3)$. At order $g$, $\Omega=1+\int\frac{dp^+}{2\pi}\int_x\big[\mf A^b_j(p^+,-\v x)\,a^{\dagger b}_j(p^+,\v x)-\mf A^{\dagger b}_j(p^+,-\v x)\,a^b_j(p^+,\v x)\big]$, and acting with $\bar\Omega^\dagger\Omega$ on the vacuum only the creation parts survive,
\begin{align}\label{eq:LOstate}
\bar\Omega^\dagger\Omega\ket0\Big|_{O(g)}=\int\frac{dp^+}{2\pi}\int_x\Big[\mf A^c_j(p^+,-\v x)-\bar{\mf A}^b_j(p^+,-\v x)\,U^{bc}(\v x)\Big]a^{\dagger c}_j(p^+,\v x)\ket0\,.
\end{align}
The first term is the gluon created by $\Omega$ before the collision, the second the gluon removed by $\bar\Omega^\dagger$, with the amplitude of the rotated charges. Projecting on $\bra{g^b_i(k^+,\v w)}$ and using $U^{ab}U^{b'b}=\delta^{ab'}$,
\begin{align}\label{eq:ALO}
\mathcal A^{a}_{{\rm LO}\,i}(k)&=\int_{w}e^{-i\v k\cdot\v w}\Big[U^{ab}(\v w)\,\mf A^b_i(k^+,-\v w)-\bar{\mf A}^a_i(k^+,-\v w)\Big]\,.
\end{align}
The leading-order spectrum in terms of the coefficients is therefore
\begin{align}\label{eq:LOcoeff}
\frac{d^3N_{\rm LO}}{d^3k}=\frac{1}{(2\pi)^3}\int_{w,w'}e^{-i\v k\cdot(\v w-\v w')}\Big\langle&\Big[\mf A^{\dagger b'}_i(k^+,-\v w')\,U^{ab'}(\v w')-\bar{\mf A}^{\dagger a}_i(k^+,-\v w')\Big]\notag\\
&\times\Big[U^{ab}(\v w)\,\mf A^b_i(k^+,-\v w)-\bar{\mf A}^a_i(k^+,-\v w)\Big]\Big\rangle\,.
\end{align}
As a check we insert the Weizs\"acker--Williams amplitude~\eqref{eq:A1X}. Then $\mathcal A^{a}_{{\rm LO}\,i}(k)=\frac{ig}{\sqrt2\,\pi\sqrt{k^+}}\int_{w,x}e^{-i\v k\cdot\v w}\,\Kv^i(\v x-\v w)\,[U(\v w)-U(\v x)]^{ab}\rho^b(\v x)$ and
\begin{align}\label{eq:LO}
\frac{d^3N_{\rm LO}}{d^3k}=\ &\frac{1}{(2\pi)^3}\,\frac{g^2}{2\pi^2}\,\frac{1}{k^+}\int_{x,x',w,w'}e^{-i\v k\cdot(\v w-\v w')}\,
\Kv(\v x-\v w)\cdot\Kv(\v x'-\v w')
\notag\\&\times
\Big\langle\big[U(\v x')-U(\v w')\big]^{ab'}\rho^{b'}(\v x')\,\big[U(\v x)-U(\v w)\big]^{ab}\rho^{b}(\v x)\Big\rangle\,.
\end{align}
The amplitude contains the difference of the Wilson lines at the position of the gluon and at the position of the charge that emits it: a gluon is produced only if the target rotates it differently from its source. For a projectile with $\langle\rho^a(\v x)\rho^b(\v y)\rangle=\delta^{ab}n(\v x,\v y)/(N_c^2-1)$, Eq.~\eqref{eq:LO} is the known leading-order result~\cite{Kovchegov:2001sc,Kovner:2006wr} with the target operator ${\rm Tr}\{[U^\dagger(\v x')-U^\dagger(\v w')][U(\v x)-U(\v w)]\}$ and the prefactor $\alpha_s/[4\pi^4(N_c^2-1)]$, $\alpha_s=g^2/4\pi$.

\subsection{Next-to-leading order: the amplitudes}\label{sec:NLOamplitudes}

\paragraph{One gluon in the final state.} At order $g^3$ we collect all terms of $\bar\Omega^\dagger\Omega$ that connect the vacuum with a one-gluon state. $\Omega\ket0$ contains the vacuum at order $1$ and $g^2$, one gluon at order $g$ and $g^3$, and two gluons at order $g^2$, Eq.~\eqref{eq:Omegavac}; $\bar\Omega^\dagger$ maps them onto a one-gluon state with $\bar{\mf A}^\dagger_1$, with $\bar{\mf N}$ and $\bar{\mf B}_3^\dagger$, and with $\bar{\mf A}_2^\dagger$ and $\bar{\mf C}_1^\dagger$, respectively. After the elimination of $\bar{\mf A}^{(3)}_1$ with Eq.~\eqref{eq:A31unit}, the terms that this relation generates combine with the direct ones, and the result has five parts,
\begin{align}\label{eq:A3amp}
\mathcal A^a_{(3)\,i}(k)=\int_w e^{-i\v k\cdot\v w}\Big\{\mathcal A_{N}+\mathcal A_{A^{(3)}}+\mathcal A_{B_3}+\mathcal A_{AB_2}+\mathcal A_{CB_2}\Big\}^a_i(\v w)\,,
\end{align}
with
\begin{align}\label{eq:A3parts}
\mathcal A_{N}&=-\bar{\mf A}^a_i(k^+,-\v w)\,\mf N_2+U^{ab}(\v w)\,\bar{\mf N}_2\,\mf A^b_i(k^+,-\v w)-\big[\bar{\mf N}_2,\bar{\mf A}^a_i(k^+,-\v w)\big]\,,
\notag\\
\mathcal A_{A^{(3)}}&=U^{ab}(\v w)\,\mf A^{(3)b}_i(k^+,-\v w)-\bar{\mf A}^{(3)a}_i(k^+,-\v w)\,,
\notag\\
\mathcal A_{B_3}&=\int\frac{dp^+}{2\pi}\int_x\bar{\mf B}^{\dagger ba}_{3ji}(k^+,\v w;p^+,-\v x)\Big[U^{bc}(\v x)\,\mf A^c_j(p^+,-\v x)-\bar{\mf A}^b_j(p^+,-\v x)\Big]\,,
\notag\\
\mathcal A_{AB_2}&=2\int\frac{dp^+}{2\pi}\int_x\Big[U^{ad}(\v w)U^{bc}(\v x)\,\bar{\mf A}^{\dagger b}_j(p^+,-\v x)\,\mf B^{cd}_{2ji}(k^+,-\v w;p^+,-\v x)
\notag\\&\hspace{2.2cm}
-\bar{\mf A}^{\dagger b}_j(p^+,-\v x)\,\bar{\mf B}^{ba}_{2ji}(k^+,-\v w;p^+,-\v x)\Big]\,,
\notag\\
\mathcal A_{CB_2}&=2\int\frac{dp^+}{2\pi}\int_{x,z}\Big[\bar{\mf C}^{\dagger dca}_{1jki}(k^+,\v w-\v z;p^+,\v z-\v x)\,U^{db}(\v x)U^{ce}(\v z)\,\mf B^{be}_{2jk}(k^+-p^+,-\v z;p^+,-\v x)
\notag\\&\hspace{2.2cm}
-\bar{\mf C}^{\dagger bca}_{1jki}(k^+,\v w-\v z;p^+,\v z-\v x)\,\bar{\mf B}^{bc}_{2jk}(k^+-p^+,-\v z;p^+,-\v x)\Big]\,.
\end{align}
$\mathcal A_N$ contains the normalisation of the state before and after the collision and the commutator that Eq.~\eqref{eq:A31unit} produces. $\mathcal A_{A^{(3)}}$ is the one-gluon component of order $g^3$ of the wave function, scattered on the target or removed after it. In $\mathcal A_{B_3}$ a gluon of momentum $p^+$ emitted at leading order is turned into the measured gluon by its interaction with the rotated valence charges after the collision. In $\mathcal A_{AB_2}$ one of the two gluons of the two-gluon component is absorbed by the valence charges after the collision, and in $\mathcal A_{CB_2}$ the two gluons merge after the collision. In each part the two terms differ by the position of the target relative to the last interaction, and their sum vanishes for $U=1$.

The commutator in $\mathcal A_N$ is the place where the order of the operators discussed in Section~\ref{sec:unitarity} enters the cross section. With $[\mf A^{(3)}]_{\rho\rho\rho}=\mf A\,\mf N_2$, Eq.~\eqref{eq:A3threerho}, and Eq.~\eqref{eq:N2Acomb} for the barred coefficients, the terms of $\mathcal A_N$ and the three-$\rho$ part of $\mathcal A_{A^{(3)}}$ add up to
\begin{align}\label{eq:Nterms}
\mathcal A_N+\big[\mathcal A_{A^{(3)}}\big]_{\rho\rho\rho}=-\bar{\mf A}\,\mf N_2+U\,\bar{\mf N}_2\,\mf A+U\,\mf A\,\mf N_2-\bar{\mf N}_2\,\bar{\mf A}\,,
\end{align}
in which no commutator is left and which vanishes for $U=1$.

\paragraph{Two gluons in the final state.} At order $g^2$ the two-gluon state is reached from the two-gluon component of $\Omega\ket0$ with $\bar{\mf N}=1$, from its one-gluon component with $\bar{\mf A}^\dagger_1$ (a second gluon is added) and $\bar{\mf C}_2^\dagger$ (the gluon splits), and from the vacuum with $\bar{\mf B}_1^\dagger$, which is eliminated with Eq.~\eqref{eq:B1unit}. The result is
\begin{align}\label{eq:A2gamp}
\mathcal A^{ab}_{ij}(k;q^+,\v v)=\frac{2}{\sqrt{2\pi}}\int_w e^{-i\v k\cdot\v w}\Big\{\mathcal A_B+\mathcal A^{(1)}_{AA}+\mathcal A^{(2)}_{AA}-\mathcal A_{CA}\Big\}^{ab}_{ij}(\v w;q^+,\v v)\,,
\end{align}
with
\begin{align}\label{eq:A2parts}
\mathcal A_B&=U^{ac}(\v w)\,\mf B^{bc}_{2ji}(k^+,-\v w;q^+,-\v v)-U^{\dagger bc}(\v v)\,\bar{\mf B}^{ca}_{2ji}(k^+,-\v w;q^+,-\v v)\,,
\notag\\
\mathcal A^{(1)}_{AA}&=\tfrac12\Big[U^{\dagger bc}(\v v)\,\bar{\mf A}^c_j(q^+,-\v v)\,\bar{\mf A}^a_i(k^+,-\v w)-U^{ac}(\v w)\,\bar{\mf A}^d_j(q^+,-\v v)\,U^{\dagger bd}(\v v)\,\mf A^c_i(k^+,-\v w)\Big]\,,
\notag\\
\mathcal A^{(2)}_{AA}&=\tfrac12\Big[U^{\dagger bc}(\v v)\,\bar{\mf A}^a_i(k^+,-\v w)\,\bar{\mf A}^c_j(q^+,-\v v)-\bar{\mf A}^a_i(k^+,-\v w)\,\mf A^b_j(q^+,-\v v)\Big]\,,
\notag\\
\mathcal A_{CA}&=\int_z\Big[\bar{\mf C}^{dac}_{1jik}(q^++k^+,\v z-\v w;q^+,\v w-\v v)\,U^{\dagger bd}(\v v)\,U^{ce}(\v z)\,\mf A^e_k(q^++k^+,-\v z)
\notag\\&\hspace{1.2cm}
-U^{\dagger bd}(\v v)\,\bar{\mf C}^{dac}_{1jik}(q^++k^+,\v z-\v w;q^+,\v w-\v v)\,\bar{\mf A}^c_k(q^++k^+,-\v z)\Big]\,,
\end{align}
where $U^{\dagger bc}=U^{cb}$. $\mathcal A_B$ is the two-gluon component of the wave function, scattered on the target or removed after it. $\mathcal A^{(1)}_{AA}$ and $\mathcal A^{(2)}_{AA}$ are products of two single emissions in the two orders that $\bar\Omega^\dagger$ generates; they do not combine into one term because the coefficients do not commute. $\mathcal A_{CA}$ is the splitting, after the collision, of a gluon of momentum $q^++k^+$ into the measured gluon and the second gluon. The colour index $b$ of the second gluon is that of the bare state; its rotation by $U(\v v)$ would drop out of the sum over final states.

We have verified Eqs.~\eqref{eq:A3amp}--\eqref{eq:A2parts} by an independent symbolic expansion (Appendix~\ref{app:rows_check}).

\subsection{The cross section at order \texorpdfstring{$g^4$}{g4}: Groups and rows}\label{sec:groups}

Inserting Eqs.~\eqref{eq:ampexp} and~\eqref{eq:A2gamp} into Eq.~\eqref{eq:N_amplitudes},
\begin{align}\label{eq:NLOdecomp}
\frac{d^3N}{d^3k}=\frac{d^3N_{\rm LO}}{d^3k}+\frac{d^3N^{\rm(I)}}{d^3k}+\frac{d^3N^{\rm(II)}}{d^3k}+\frac{d^3N^{\rm(III)}}{d^3k}+O(g^6)\,,
\end{align}
where
\begin{align}\label{eq:Groups}
\frac{d^3N^{\rm(I)}}{d^3k}&=\frac{1}{(2\pi)^3}\Big\langle\mathcal A^{a\dagger}_{{\rm LO}\,i}(k)\,\mathcal A^a_{(3)\,i}(k)\Big\rangle+{\rm c.c.}\,,\notag\\
\frac{d^3N^{\rm(II)}}{d^3k}&=\frac{1}{(2\pi)^3}\int dp^+\!\int_z\Big\langle\big|\mathcal A_B\big|^2+\big|\mathcal A^{(1)}_{AA}\big|^2+\big|\mathcal A^{(2)}_{AA}\big|^2+\big|\mathcal A_{CA}\big|^2\Big\rangle\,,\notag\\
\frac{d^3N^{\rm(III)}}{d^3k}&=\frac{1}{(2\pi)^3}\int dp^+\!\int_z\Big\langle\mathcal A^{(1)\dagger}_{AA}\mathcal A_B+\mathcal A^{(2)\dagger}_{AA}\mathcal A_B-\mathcal A^{\dagger}_{CA}\mathcal A_B
\notag\\&\hspace{3.3cm}+\mathcal A^{(2)\dagger}_{AA}\mathcal A^{(1)}_{AA}-\mathcal A^{\dagger}_{CA}\mathcal A^{(1)}_{AA}-\mathcal A^{\dagger}_{CA}\mathcal A^{(2)}_{AA}\Big\rangle+{\rm c.c.}
\end{align}
Here the momentum and the position of the second gluon have been renamed, $(q^+,\v v)\to(p^+,\v z)$, and the products of the parts of the two-gluon amplitude stand for
\begin{align}\label{eq:XYdef}
\mathcal X^\dagger\mathcal Y\equiv\frac{2}{\pi}\int_{w,w'}e^{-i\v k\cdot(\v w-\v w')}\,\Big[\mathcal X^{ab}_{ij}(\v w';p^+,\v z)\Big]^\dagger\,\mathcal Y^{ab}_{ij}(\v w;p^+,\v z)\,,\qquad|\mathcal X|^2\equiv\mathcal X^\dagger\mathcal X\,,
\end{align}
where the factor $2/\pi$ is the square of the prefactor of Eq.~\eqref{eq:A2gamp}. Group~I is the interference of the leading-order amplitude with the one-gluon amplitude of order $g^3$; the second gluon is virtual. Groups~II and~III have two gluons in the final state: Group~II contains the squares of the four parts of the two-gluon amplitude, Group~III their six interferences. Groups~I and~III are written with their complex conjugates, i.e.\ the Hermitian conjugate of the operator inside the average; Group~II is Hermitian by itself. The square of $\mathcal A_{(3)}$ is of order $g^6$.

We call the individual products of coefficients \emph{rows}. Group~I has five rows, one for each part of Eq.~\eqref{eq:A3parts}; Group~II has four and Group~III six. They are listed in Table~\ref{tab:rows}. As a representative of each group we display one row in terms of the coefficients; the complete list is given in Appendix~\ref{app:rows}.

\paragraph{Group I, Row I.III.} The interference of the leading-order amplitude with $\mathcal A_{B_3}$ is
\begin{align}\label{eq:rowI_III}
\frac{d^3N^{\rm(I.III)}}{d^3k}=\ &\frac{1}{(2\pi)^3}\int_{w,w'}e^{-i\v k\cdot(\v w-\v w')}\int\frac{dp^+}{2\pi}\int_x\Big\langle\Big[\mf A^{\dagger b'}_i(k^+,-\v w')\,U^{ab'}(\v w')-\bar{\mf A}^{\dagger a}_i(k^+,-\v w')\Big]\notag\\
&\times\bar{\mf B}^{\dagger ba}_{3ji}(k^+,\v w;p^+,-\v x)\Big[U^{bc}(\v x)\,\mf A^c_j(p^+,-\v x)-\bar{\mf A}^b_j(p^+,-\v x)\Big]\Big\rangle+{\rm c.c.}
\end{align}

\paragraph{Group II, Row II.I.} The square of $\mathcal A_B$ is
\begin{align}\label{eq:rowII_I}
\frac{d^3N^{\rm(II.I)}}{d^3k}=\ &\frac{4}{(2\pi)^3}\int\frac{dp^+}{2\pi}\int_{z,w,w'}e^{-i\v k\cdot(\v w-\v w')}\notag\\
&\times\Big\langle\Big[\mf B^{\dagger bc'}_{2ji}(k^+,-\v w';p^+,-\v z)\,U^{ac'}(\v w')-\bar{\mf B}^{\dagger c'a}_{2ji}(k^+,-\v w';p^+,-\v z)\,U^{c'b}(\v z)\Big]\notag\\
&\quad\times\Big[U^{ac}(\v w)\,\mf B^{bc}_{2ji}(k^+,-\v w;p^+,-\v z)-U^{cb}(\v z)\,\bar{\mf B}^{ca}_{2ji}(k^+,-\v w;p^+,-\v z)\Big]\Big\rangle\,.
\end{align}

\paragraph{Group III, Row III.I.} The interference of $\mathcal A^{(1)}_{AA}$ with $\mathcal A_B$ is
\begin{align}\label{eq:rowIII_I}
\frac{d^3N^{\rm(III.I)}}{d^3k}=\ &\frac{2}{(2\pi)^3}\int\frac{dp^+}{2\pi}\int_{z,w,w'}e^{-i\v k\cdot(\v w-\v w')}\notag\\
&\times\Big\langle\Big[\bar{\mf A}^{\dagger a}_i(k^+,-\v w')\,\bar{\mf A}^{\dagger c'}_j(p^+,-\v z)\,U^{c'b}(\v z)-\mf A^{\dagger c'}_i(k^+,-\v w')\,U^{d'b}(\v z)\,\bar{\mf A}^{\dagger d'}_j(p^+,-\v z)\,U^{ac'}(\v w')\Big]\notag\\
&\quad\times\Big[U^{ac}(\v w)\,\mf B^{bc}_{2ji}(k^+,-\v w;p^+,-\v z)-U^{cb}(\v z)\,\bar{\mf B}^{ca}_{2ji}(k^+,-\v w;p^+,-\v z)\Big]\Big\rangle+{\rm c.c.}
\end{align}
In these expressions the conjugated factor stands to the left, at the primed point, and the order of the coefficients is the order in which they act.

\begin{table}[t]
\centering
\renewcommand{\arraystretch}{1.25}
\footnotesize
\begin{tabular}{@{}l l l c l l@{}}
\toprule
Row & product & coefficients & $\#\rho$ & type & range of $p^+$\\
\midrule
I.I & $\mathcal A_{\rm LO}^\dagger\times$ Eq.~\eqref{eq:Nterms} & $\mf A^\dagger\,\mf A\,\mf N_2$ & 4 & virtual & ---\\
I.II & $\mathcal A_{\rm LO}^\dagger\,\mathcal A_{A^{(3)}}$ (one, two $\rho$) & $\mf A^\dagger\,\mf A^{(3)}$ & 2, 3 & virtual & ---\\
I.III & $\mathcal A_{\rm LO}^\dagger\,\mathcal A_{B_3}$ & $\mf A^\dagger\,\bar{\mf B}_3^\dagger\,\mf A$ & 3, 4 & virtual & $\Lambda<p^+<\vee$, $|p^+-k^+|>\Lambda$\\
I.IV & $\mathcal A_{\rm LO}^\dagger\,\mathcal A_{AB_2}$ & $\mf A^\dagger\,\bar{\mf A}^\dagger\,\mf B_2$ & 3, 4 & virtual & $\Lambda<p^+<\vee-k^+$ $[\vee]$\\
I.V & $\mathcal A_{\rm LO}^\dagger\,\mathcal A_{CB_2}$ & $\mf A^\dagger\,\bar{\mf C}_1^\dagger\,\mf B_2$ & 2, 3 & virtual & $\Lambda<p^+<k^+-\Lambda$\\
\midrule
II.I & $|\mathcal A_B|^2$ & $\mf B_2^\dagger\,\mf B_2$ & 2--4 & real & $\Lambda<p^+<\vee-k^+$ $[\vee]$\\
II.II & $|\mathcal A^{(1)}_{AA}|^2$ & $\mf A^\dagger\mf A^\dagger\,\mf A\,\mf A$ & 4 & real & $\Lambda<p^+<\vee$\\
II.III & $|\mathcal A^{(2)}_{AA}|^2$ & $\mf A^\dagger\mf A^\dagger\,\mf A\,\mf A$ & 4 & real & $\Lambda<p^+<\vee$\\
II.IV & $|\mathcal A_{CA}|^2$ & $\mf A^\dagger\,\bar{\mf C}_1^\dagger\,\bar{\mf C}_1\,\mf A$ & 2 & real & $\Lambda<p^+<\vee-k^+$\\
\midrule
III.I & $\mathcal A^{(1)\dagger}_{AA}\,\mathcal A_B$ & $\mf A^\dagger\mf A^\dagger\,\mf B_2$ & 3, 4 & real & $\Lambda<p^+<\vee-k^+$ $[\vee]$\\
III.II & $\mathcal A^{(2)\dagger}_{AA}\,\mathcal A_B$ & $\mf A^\dagger\mf A^\dagger\,\mf B_2$ & 3, 4 & real & $\Lambda<p^+<\vee-k^+$ $[\vee]$\\
III.III & $-\mathcal A^{\dagger}_{CA}\,\mathcal A_B$ & $\mf A^\dagger\,\bar{\mf C}_1^\dagger\,\mf B_2$ & 2, 3 & real & $\Lambda<p^+<\vee-k^+$\\
III.IV & $\mathcal A^{(2)\dagger}_{AA}\,\mathcal A^{(1)}_{AA}$ & $\mf A^\dagger\mf A^\dagger\,\mf A\,\mf A$ & 4 & real & $\Lambda<p^+<\vee$\\
III.V & $-\mathcal A^{\dagger}_{CA}\,\mathcal A^{(1)}_{AA}$ & $\mf A^\dagger\,\bar{\mf C}_1^\dagger\,\mf A\,\mf A$ & 3 & real & $\Lambda<p^+<\vee-k^+$\\
III.VI & $-\mathcal A^{\dagger}_{CA}\,\mathcal A^{(2)}_{AA}$ & $\mf A^\dagger\,\bar{\mf C}_1^\dagger\,\mf A\,\mf A$ & 3 & real & $\Lambda<p^+<\vee-k^+$\\
\bottomrule
\end{tabular}
\caption{The rows of the cross section at order $g^4$. ``Coefficients'' lists the types of coefficients in the product, barred or unbarred. $\#\rho$ is the number of charges before any reordering, with $\mf B_2$ and $\mf B_3$ containing parts with one and with two charges. ``Virtual'': one gluon in the final state; ``real'': two gluons in the final state, the second integrated over. $p^+$ is the longitudinal momentum of the second gluon; where two ranges are given, the first applies to the part of $\mf B_2$ linear in $\rho$ and the bracket $[\vee]$, i.e.\ $\Lambda<p^+<\vee$, to its $\{\rho,\rho\}$ part. Groups~I and~III are to be taken with the complex conjugate. The three-$\rho$ part of $\mathcal A_{A^{(3)}}$ is included in Row~I.I.}
\label{tab:rows}
\end{table}

\paragraph{Number of charges and the four-$\rho$ terms.} The number of charges in a row follows from the coefficients: $\mf A$ and the linear parts of $\mf B_2$, $\mf B_3$ contain one charge, $\mf N_2$ and the $\{\rho,\rho\}$ parts of $\mf B_2$, $\mf B_3$ two, $\mf C_1$ none, and $\mf A^{(3)}$ one, two or three. Terms with four charges, the maximal number at order $g^4$, occur in ten rows: I.I, I.III, I.IV, II.I, II.II, II.III, III.I, III.II and III.IV, and in Row~I.II through the three-$\rho$ part of $\mf A^{(3)}$, which we have included in Row~I.I, Eq.~\eqref{eq:Nterms}. Rows~I.I, II.II, II.III and III.IV consist of four-$\rho$ terms only. Rows~I.V, II.IV and III.III are the only ones with terms quadratic in $\rho$ besides I.II and II.I; Rows~III.V and III.VI are cubic. The ranges of the longitudinal integrals in Table~\ref{tab:rows} follow from the supports of the coefficients, i.e.\ from the kinematic restrictions of Section~\ref{sec:lcpt}.

\NOTE{The assignment of the three-$\rho$ part of $\mathcal A_{A^{(3)}}$ to Row~I.I follows the current version of the paper; the calculation file keeps it in the second row. The sum is the same.}

Eqs.~\eqref{eq:Groups}, with the amplitudes~\eqref{eq:A3parts} and~\eqref{eq:A2parts} and the coefficients of Table~\ref{tab:coeffs}, are the complete cross section at order $g^4$ in terms of the coefficients of $\Omega$. Its evaluation is the subject of the following sections.
